% ============================================================================================================
% DRAFTS for the parts of the report that still hold template placeholder text: title, abstract, introduction
% (with a notation paragraph and a roadmap), and conclusion. Written at ISL level for the team to edit; every
% number is from results/ (Part 1 and the data summaries). The error-bar work is mentioned only as upcoming.
% Section numbers assume the draft's order: 2 Background, 3 Data Cleaning, 4 Technical Implementation, 5 Results.
% ============================================================================================================

% ---- Title (replaces "Your Title Should Go Here")
% \title{Can We Predict a Match Before Two Strangers Meet? Evidence from a Speed-Dating Experiment}
% \author{Sophia Qin, Jenna Yang, and Abraham Yeung}

% ---- Abstract (replaces the template text inside \begin{abstract} ... \end{abstract})
Dating apps and matchmakers try to predict which pairs of people will like each other before they meet. Using the
Columbia speed-dating experiment (21 events, 551 participants and 4,184 woman--man pairs), we ask how well a
\emph{mutual} match can be predicted from information available before the date, and how much the four-minute date
adds. We compare logistic regression, stepwise selection, the lasso, linear and quadratic discriminant analysis and
random forests, and we hold out whole events in cross-validation so that every prediction is for people the model has
never seen. Before the date the best model reaches an AUC of 0.58, only slightly better than chance; adding the
ratings the two participants gave each other after the date raises it to 0.85. Features designed to measure
compatibility, such as similar interests or a small age gap, add about 0.01. Participants' choices depended far more
on attractiveness than their stated preferences suggested. Every hypothesis was registered before it was tested.

% ---- Introduction (replaces the template text under \section{Introduction})
\section{Introduction}

Online dating platforms match millions of people using profiles: age, interests, what each person says they want.
Behind every recommendation is a prediction that two strangers will like each other. How good can such predictions
be? Speed dating offers an unusually clean test. In the Columbia experiment \cite[Section II]{Fisman:2006}, every
participant filled in a detailed questionnaire, then met each partner of the opposite sex for four minutes and
privately decided whether to see that person again. A \emph{match} occurs when both say yes, so for every pair we
know what was known beforehand, what each person thought after meeting, and whether the interest was mutual.

Our main question is how well a mutual match can be predicted from information available before the date, and which
pair characteristics matter. A second question is how much the four-minute meeting adds. To answer them we built two
sets of predictors, one known before the date and one that adds the participants' ratings of each other after it,
and fitted the classification methods of the course textbook to each. The data need one departure from textbook
practice: the same person appears in up to 22 pairs, so we hold out whole events, not individual rows, when we
measure accuracy.

We find that pre-date information carries a real but small signal. The best model reaches an AUC of 0.58, where 0.5 is
a coin flip, and explicitly modelling compatibility barely helps. The four-minute meeting changes the picture: with
the in-date ratings the best AUC is 0.85. Consistent with earlier work \cite{Fisman:2006,Eastwick:2008}, attractiveness
drives decisions far more than people say it does. Together these results suggest that compatibility emerges from the
interaction itself rather than from characteristics a profile can capture, which echoes Joel, Eastwick and Finkel's
finding that pairwise attraction is hard to predict from pre-date reports \cite{Joel:2017}.

\emph{Notation and terminology.} A \emph{dyad} is one woman--man pair who met at an event; there are $n=4{,}184$ dyads in
$G=21$ events (also called waves). For dyad $i$, $Y_i=1$ if both people said yes (a match) and $0$ otherwise, and $X_i$
is its vector of predictors. Variables of the woman carry the prefix \texttt{f\_} and those of the man \texttt{m\_}.
The feature set \texttt{pre} contains only information available before the date; \texttt{pre+date} adds the ratings
the two people gave each other after it. All accuracy measures are computed on held-out events, and AUC denotes the
area under the ROC curve.

\emph{Roadmap.} Section~2 describes the data and reviews the statistical methods we use, including a few that go
beyond the course textbook. Section~3 explains how we cleaned the data and built one row per dyad. Section~4 describes
the predictors, the cross-validation design and the models. Section~5 reports the results and how robust they are,
and Section~6 concludes and outlines the second half of the project.

% ---- Conclusion (replaces the template text under \section{Conclusion})
\section{Conclusion}

We asked how well a mutual match between two strangers can be predicted before they meet. Using the Columbia
speed-dating data and holding out whole events, so that every prediction concerns people the model has never seen,
we found that pre-date information gives an AUC of about 0.58 at best, compatibility features add about 0.01, and the
ratings from a four-minute date raise the AUC to 0.85. People's actual choices weighted attractiveness roughly twice as
heavily as their stated preferences. For matchmaking systems the lesson is that profiles say a little about who will
be liked, but much more is learned once two people have met.

These conclusions come with limits. The data are one experiment, at one university, with 21 events. The best
pre-date model clears our pre-registered threshold only narrowly, and its margin over logistic regression depends
on how missing data are handled (Section~5). In the second half of the project we will study how much to trust
cross-validated results when the same people appear in many rows, and we will extend the analysis to
participants' own predictions and to what a matchmaker using our models could gain.
